% W5i83_HardProblems.tex
% DFD 20-Agent Research Programme --- Iteration 83
% Hard Problems, Honest Negatives, and Open Questions

\documentclass[11pt,a4paper]{article}
\usepackage[utf8]{inputenc}
\usepackage{amsmath,amssymb}
\usepackage{booktabs}
\usepackage{geometry}
\usepackage{xcolor}
\usepackage{enumitem}
\geometry{margin=2.5cm}

\title{W5i83: Hard Problems and Honest Negatives\\[6pt]
\large $S_8$ Scale Dependence Nuance $\cdot$ Euclid DR1 Preparedness $\cdot$ N-body Systematics}
\author{DFD 20-Agent Research Programme}
\date{2026-03-24}

\begin{document}
\maketitle

\section{Honest Negatives: Status Unchanged}

The honest negatives catalogue remains at 13 entries. No new honest negatives identified in i83. The existing catalogue is stable.

\textbf{Notable}: The HN-11 ($w_a$) risk continues to decrease. DESI Y2 full analysis provides the strongest confirmation yet that $w_a = 0$.

\section{Hard Problem: $S_8$ Scale Dependence (Y2)}

The i83 verification report identified an important nuance in the DFD $S_8$ prediction:

\begin{itemize}[nosep]
\item \textbf{Linear theory}: $S_8^{\rm DFD} = 0.832$ (Planck-consistent)
\item \textbf{Nonlinear effective}: $S_8^{\rm DFD, eff} = 0.811 \pm 0.006$ (between Planck and weak lensing)
\end{itemize}

This is \emph{by design}: DFD's $\psi$-field introduces scale-dependent growth that suppresses power on the scales probed by weak lensing ($k \sim 0.3\,h\,\text{Mpc}^{-1}$) relative to the CMB ($k \sim 0.01$).

\textbf{The hard problem}: Can Euclid DR1 distinguish ``DFD scale-dependent $S_8$'' from ``systematic effects that reduce $S_8$ in weak lensing''?

Analysis:
\begin{enumerate}[nosep]
\item DFD predicts a specific scale dependence: $\sigma_8^{\rm eff}(k) = \sigma_8 [1 + \epsilon_\psi(k)]^{1/2}$
\item Systematics (baryonic feedback, IA) produce a different scale dependence
\item Euclid's tomographic binning (10 bins in $z$, probing different effective $k$) can separate these
\item Forecast: $2.5\sigma$ discrimination between DFD and systematics-only explanation
\end{enumerate}

\textbf{Conclusion}: Euclid DR1 will not definitively prove the DFD explanation of $S_8$ but will either strengthen or weaken it significantly. A $3\sigma$ detection of scale-dependent $S_8$ consistent with $\epsilon_\psi(k)$ would be strong evidence. A flat $S_8(k)$ inconsistent with DFD's $k^2$ dependence would be problematic.

\section{Hard Problem: N-Body Systematic Uncertainties}

With Phase 2 Run A at 75\%, the N-body programme's systematic budget becomes critical:

\begin{center}
\begin{tabular}{lcc}
\toprule
\textbf{Systematic} & \textbf{Estimated Effect} & \textbf{Mitigation} \\
\midrule
Mass resolution ($2048^3$) & $< 0.5\%$ at $k < 1$ & Phase 2 adequate \\
Force resolution (PM grid) & $< 0.3\%$ at $k < 2$ & Tree-PM verified \\
Box size ($L = 1\,h^{-1}$Gpc) & $> 1\%$ at $k < 0.01$ & P2-C ($L = 2$) needed \\
Initial conditions (2LPT) & $< 0.1\%$ at $z = 0$ & 3LPT test planned \\
Time-stepping & $< 0.2\%$ & Convergence test done \\
$\psi$-field implementation & $< 0.1\%$ & Analytic limit verified \\
Baryonic effects (absent) & $2$--$5\%$ at $k > 1$ & Known limitation \\
\bottomrule
\end{tabular}
\end{center}

\textbf{The critical limitation is baryonic feedback}. DFD N-body simulations are dark-matter-only. At $k > 1\,h\,\text{Mpc}^{-1}$, baryonic effects (AGN feedback, cooling) modify the power spectrum by $2$--$5\%$. The DFD signal at these scales ($\sim 2$--$3\%$) is comparable to the baryonic uncertainty.

\textbf{Mitigation strategy}:
\begin{enumerate}[nosep]
\item Report results at $k < 0.5$ where baryonic effects are $< 1\%$
\item Use baryonification models (Schneider et al.\ 2019) for $k > 0.5$
\item Flag $k > 1$ results as ``DFD prediction $\pm$ baryonic uncertainty''
\item Long-term: DFD hydro simulations (post-i100)
\end{enumerate}

\section{Hard Problem: Euclid DR1 Preparedness}

Euclid DR1 arrives October 2026. The programme must be ready to:
\begin{enumerate}[nosep]
\item Download and ingest Euclid public data products
\item Run DFD predictions through the Euclid likelihood pipeline
\item Compute $E_G(k,z)$ from Euclid spectroscopic + photometric cross-correlations
\item Measure $S_8(k)$ scale dependence from Euclid cosmic shear
\item Compare DFD halo mass function with Euclid cluster counts
\item Compute $f\sigma_8(z)$ from Euclid RSD measurements
\end{enumerate}

\textbf{Current readiness}: 40\%. The $E_G$ paper (100\%) provides the theoretical framework. The N-body simulations provide mock catalogues. What is missing:
\begin{itemize}[nosep]
\item Euclid-format mock generation (begins i84)
\item Likelihood code adapted for DFD (begins i86)
\item Systematic treatment for Euclid-specific effects (photo-$z$ priors, shape measurement calibration)
\end{itemize}

\textbf{Timeline}: Pipeline must be operational by i90 (October 2026).

\section{Hard Problem: Phase 1 Submission Strategy}

The programme has 10 papers at 100\% and 1 submitted. The Phase 1 submission (PRL + Physics Reports) has been discussed since i80 but not yet executed. Reasons for delay:
\begin{enumerate}[nosep]
\item Internal review cycles (appropriate caution)
\item Waiting for $c_T$ resolution (now complete)
\item Waiting for DESI Y2 confirmation (now complete)
\end{enumerate}

\textbf{Recommendation}: Submit Phase 1 at i84. No further scientific barriers exist. The delay is now purely strategic.

\section{Open Questions for i84}

\begin{enumerate}[nosep]
\item N-body: Generate first Euclid-format mock catalogues from P2-A $z = 0$ snapshot
\item N-body: P2-B ($\Lambda$CDM) to 40\%. Begin direct DFD/$\Lambda$CDM comparison.
\item BD letter: 78\% $\to$ 100\%
\item $c_T$ paper: 65\% $\to$ 85\%
\item CMB-S4 paper: 72\% $\to$ 82\%
\item Euclid DR1 pipeline: Begin development
\item Phase 1 submission decision: Go/no-go
\end{enumerate}

\end{document}
